%% text-sharing.tex -- shared closing section
\subsubsection{Building and sharing}

The package is \texttt{florilegium.sty} with twelve Lua files beside it. Of these,
\texttt{florilegium.lua} is the glue to \TeX. The names of the other eleven begin with
\texttt{florilegium-}, because \TeX\ finds a file by its bare name, the document's folder
first, and a generic name would lose to a file of the author's own. After the prefix come
\texttt{ink} (geometry, color and PDF output), \texttt{flora} (the hand-designed primitives),
\texttt{species} (the genomes), \texttt{native} (the species' organs), \texttt{pgfmath} (the
few numbers taken from \TeX's arithmetic), \texttt{authored} and \texttt{grow} (the two
compositions), \texttt{initial} (the decorated initials), and two themes with the figures one
of them draws from (\texttt{theme-geometric}, the same plant cut from flat shapes;
\texttt{theme-math}, figures, symbols and equations cascading down the margin;
\texttt{figures}). Two more files, \texttt{tikzlibraryflorilegium.code.tex} and its tables,
are the TikZ library, which draws the species in pure TikZ. The package needs Lua\LaTeX\ and
one compilation, at about a second per page. To share it, the three routes that work are, in
increasing order of effort: paste the files into a blank Overleaf project with the compiler set
to Lua\LaTeX\ and share the link; put them in a Git repository with a continuous integration
job that compiles the PDFs on every push and attaches them to tagged releases; or publish the
package so that any document can \texttt{\textbackslash usepackage} the garden and choose its
own seed, column and hue.

\begin{remark}
Three mistakes cost time and are worth recording. A literal written into the page bypasses
TikZ, so it also bypasses TikZ's transformations: a primitive called inside
\verb|\begin{scope}[rotate=30]| must read pgf's current matrix itself (it is in
\verb|\pgf@pt@aa| to \verb|\pgf@pt@bb| and the dimensions \verb|\pgf@pt@x|,
\verb|\pgf@pt@y|) and apply it in Lua. A Lua function that returns two values, placed
in the middle of a table constructor such as \verb|{x0, y0, f(t), x1, y1}|, contributes only
its first value; the rest of the table shifts, and the error surfaces three functions
later as arithmetic on \texttt{nil}. And opening a scope in pgf empties the current path, so a
path built just before the scope and filled inside it never reaches the page. The TikZ library
drew every leaf without its darker half that way, and nobody noticed until its drawings were
compared, in total darkness, with the Lua's, which keeps the path.
\end{remark}
